% GENERATED FILE. Edit canonical lessons, then run npm run generate:books.
% canonical-source-hash: fnv1a64:e4318a48ce798622
% canonical-lessons: LA-C227-fatum, LA-C227-cubile, LA-C227-hospitium, LA-C227-deversorium, LA-C227-viator

\chapter{Fate, Lair, Hospitality, Lodging house, Traveller}
\label{ch:la-fate-lair-hospitality-lodging-house-traveller}
\hlchaptermodality{227}

\begin{chapteropening}
\textbf{By the end of this chapter:} \emph{I can say fate, a lair, hospitality, a lodging house and a traveller.}

Complete the last lesson of chapter 227: I can say fate, a lair, hospitality, a lodging house and a traveller.
\end{chapteropening}

\section[fātum, fātī]{fātum, fātī — fate}

\label{lesson:LA-C227-fatum}



\begin{quote}
Before the new one: say the Latin for a question, then the Latin for an answer.
\end{quote}

\subsection*{You'll want to know: fātum, fātī}
\textbf{fātum, fātī} — \textquotedblleft{}fate\textquotedblright{}.

\begin{cousinweb}
\textbf{fātum} is the past participle of \emph{fārī}, \textquotedblleft{}to speak\textquotedblright{} (EtymDB 2.1), used as a noun: what has been spoken, the word of the gods.
\end{cousinweb}

\begin{grammarlens}[title={What you've built}]
One more word for this chapter.
\end{grammarlens}

\subsection*{Your turn}
\begin{itemize}
\raggedright
  \item \emph{Say it:} \emph{fātum, fātī}
  \item \emph{Say it:} \emph{fātum, fātī}, once more
  \item \emph{Say it:} say \emph{fātum, fātī}
  \item \emph{Recall:} say the Latin for a question, then the Latin for an answer, then say \emph{fātum, fātī} again
\end{itemize}

\begin{tcolorbox}[breakable,colback=teal!4,colframe=teal!35!black,title={Before you move on}]
What does fātum, fātī mean? (\textquotedblleft{}Fate\textquotedblright{}.) Say it once more.
\end{tcolorbox}
\section[cubīle, cubīlis]{cubīle, cubīlis — a lair, a resting place}

\label{lesson:LA-C227-cubile}



\begin{quote}
Before the new one: say the Latin for an answer, then the Latin for fate.
\end{quote}

\subsection*{You'll want to know: cubīle, cubīlis}
\textbf{cubīle, cubīlis} — \textquotedblleft{}a lair, a resting place\textquotedblright{}.

\begin{cousinweb}
\textbf{cubīle} is built on \emph{cubō}, \textquotedblleft{}to lie down\textquotedblright{} (Lewis \& Short), with \emph{-īle}, an ending for a place: a place to lie.
\end{cousinweb}

\begin{grammarlens}[title={What you've built}]
One more word for this chapter.
\end{grammarlens}

\subsection*{Your turn}
\begin{itemize}
\raggedright
  \item \emph{Say it:} \emph{cubīle, cubīlis}
  \item \emph{Say it:} \emph{cubīle, cubīlis}, once more
  \item \emph{Say it:} say \emph{cubīle, cubīlis}
  \item \emph{Recall:} say the Latin for an answer, then the Latin for fate, then say \emph{cubīle, cubīlis} again
\end{itemize}

\begin{tcolorbox}[breakable,colback=teal!4,colframe=teal!35!black,title={Before you move on}]
What does cubīle, cubīlis mean? (\textquotedblleft{}A lair, a resting place\textquotedblright{}.) Say it once more.
\end{tcolorbox}
\section[hospitium, hospitiī]{hospitium, hospitiī — hospitality; lodging}

\label{lesson:LA-C227-hospitium}



\begin{quote}
Before the new one: say the Latin for fate, then the Latin for a lair.
\end{quote}

\subsection*{You'll want to know: hospitium, hospitiī}
\textbf{hospitium, hospitiī} — \textquotedblleft{}hospitality; lodging\textquotedblright{}.

\begin{cousinweb}
\textbf{hospitium} is built on \emph{hospes}, \emph{hospitis}, \textquotedblleft{}a guest, a host\textquotedblright{} (Lewis \& Short), with \emph{-ium}, an ending for a state or place.
\end{cousinweb}

\begin{grammarlens}[title={What you've built}]
One more word for this chapter.
\end{grammarlens}

\subsection*{Your turn}
\begin{itemize}
\raggedright
  \item \emph{Say it:} \emph{hospitium, hospitiī}
  \item \emph{Say it:} \emph{hospitium, hospitiī}, once more
  \item \emph{Say it:} say \emph{hospitium, hospitiī}
  \item \emph{Recall:} say the Latin for fate, then the Latin for a lair, then say \emph{hospitium, hospitiī} again
\end{itemize}

\begin{tcolorbox}[breakable,colback=teal!4,colframe=teal!35!black,title={Before you move on}]
What does hospitium, hospitiī mean? (\textquotedblleft{}Hospitality; lodging\textquotedblright{}.) Say it once more.
\end{tcolorbox}
\section[dēversōrium, dēversōriī]{dēversōrium, dēversōriī — a lodging house}

\label{lesson:LA-C227-deversorium}



\begin{quote}
Before the new one: say the Latin for a lair, then the Latin for hospitality; lodging.
\end{quote}

\subsection*{You'll want to know: dēversōrium, dēversōriī}
\textbf{dēversōrium, dēversōriī} — \textquotedblleft{}a lodging house\textquotedblright{}.

\begin{cousinweb}
\textbf{dēversōrium} is built on \emph{dēvertō}, \textquotedblleft{}to turn aside\textquotedblright{} (EtymDB 2.1), from \emph{dē}, \textquotedblleft{}away\textquotedblright{}, and \emph{vertō}, \textquotedblleft{}to turn\textquotedblright{}: a place where travellers turn aside for the night.
\end{cousinweb}

\begin{grammarlens}[title={What you've built}]
One more word for this chapter.
\end{grammarlens}

\subsection*{Your turn}
\begin{itemize}
\raggedright
  \item \emph{Say it:} \emph{dēversōrium, dēversōriī}
  \item \emph{Say it:} \emph{dēversōrium, dēversōriī}, once more
  \item \emph{Say it:} say \emph{dēversōrium, dēversōriī}
  \item \emph{Recall:} say the Latin for a lair, then the Latin for hospitality; lodging, then say \emph{dēversōrium, dēversōriī} again
\end{itemize}

\begin{tcolorbox}[breakable,colback=teal!4,colframe=teal!35!black,title={Before you move on}]
What does dēversōrium, dēversōriī mean? (\textquotedblleft{}A lodging house\textquotedblright{}.) Say it once more.
\end{tcolorbox}
\section[viātor, viātōris]{viātor, viātōris — a traveller}

\label{lesson:LA-C227-viator}



\begin{quote}
Before the new one: say the Latin for hospitality; lodging, then the Latin for a lodging house.
\end{quote}

\subsection*{You'll want to know: viātor, viātōris}
\textbf{viātor, viātōris} — \textquotedblleft{}a traveller\textquotedblright{}.

\begin{cousinweb}
\textbf{viātor} is built on \emph{viō}, \textquotedblleft{}to travel\textquotedblright{} (Lewis \& Short), from \emph{via}, \textquotedblleft{}a road\textquotedblright{}, with \emph{-tor}, \textquotedblleft{}one who does it\textquotedblright{}. \emph{via} is usually traced to Proto-Indo-European \emph{*wegh-}, \textquotedblleft{}to go, to carry\textquotedblright{} (EtymDB 2.1).
\end{cousinweb}

\begin{grammarlens}[title={What you've built}]
One more word for this chapter.
\end{grammarlens}

\subsection*{Your turn}
\begin{itemize}
\raggedright
  \item \emph{Say it:} \emph{viātor, viātōris}
  \item \emph{Say it:} \emph{viātor, viātōris}, once more
  \item \emph{Say it:} say \emph{viātor, viātōris}
  \item \emph{Recall:} say the Latin for hospitality; lodging, then the Latin for a lodging house, then say \emph{viātor, viātōris} again
\end{itemize}

\begin{tcolorbox}[breakable,colback=teal!4,colframe=teal!35!black,title={Before you move on}]
What does viātor, viātōris mean? (\textquotedblleft{}A traveller\textquotedblright{}.) Say it once more.
\end{tcolorbox}
